\chapter{Mekanisme Reaksi: Panah Lengkung}\label{ch:g12:curly-arrows}

\omterm{def:g8:balanced-equations:equation}{Persamaan setara} itu ibarat foto pertama dan foto terakhir sebuah film:
persamaan itu memperlihatkan \omterm{def:g7:chemical-reactions:reactant}{pereaksi} sebelumnya dan produk sesudahnya,
tetapi tidak memperlihatkan apa pun di antaranya. Padahal \omterm{def:g7:atoms-and-molecules:molecule}{molekul} tidak
sekadar bertukar \omterm{def:g7:atoms-and-molecules:atom}{atom}. Ikatan putus dan terbentuk satu per satu, ketika
pasangan-pasangan \omterm{def:g9:inside-the-atom:nucleus}{elektron} berpindah dari satu tempat ke tempat lain,
melalui spesies berumur pendek yang tidak pernah muncul dalam
persamaan. Ahli kimia menggambarkan perpindahan itu dengan \omterm{def:g12:curly-arrows:curly-arrow}{panah
lengkung}. Dengan beberapa aturan, panah-panah ini menjelaskan mengapa
sebuah reaksi terjadi di tempat ia terjadi, dan meramalkan reaksi yang
belum pernah diamati.

\begin{recall}
\omterm{def:g11:polarity-and-cohesion:electronegativity}{Keelektronegatifan}, \omterm{def:g11:polarity-and-cohesion:polar-bond}{ikatan polar}, dan \omterm{def:g11:polarity-and-cohesion:polar-bond}{muatan parsial}
(\cref{ch:g11:polarity-and-cohesion}). \omterm{def:g10:lewis-and-shape:lewis-structure}{Struktur Lewis}: \omterm{def:g10:lewis-and-shape:lone-pair}{pasangan
elektron ikatan} dan \omterm{def:g10:lewis-and-shape:lone-pair}{pasangan elektron bebas}
(\cref{ch:g10:lewis-and-shape}). \omterm{def:g11:functional-groups:functional-group}{Gugus fungsi} dan keluarga senyawa
(\cref{ch:g11:functional-groups}).
\end{recall}

\section{Situs kaya elektron dan situs miskin elektron}

\begin{definition}[Situs donor elektron dan situs akseptor elektron]\label{def:g12:curly-arrows:site}
Sebuah \emph{situs donor elektron}\index{situs donor elektron} pada
sebuah \omterm{def:g7:atoms-and-molecules:molecule}{molekul} atau \omterm{def:g9:ions:ion}{ion} adalah tempat yang kaya \omterm{def:g9:inside-the-atom:nucleus}{elektron}, yang mampu
memberikan sepasang \omterm{def:g9:inside-the-atom:nucleus}{elektron} untuk membentuk ikatan baru: \omterm{def:g10:lewis-and-shape:lone-pair}{pasangan
elektron bebas}, \omterm{def:g10:lewis-and-shape:multiple-bond}{ikatan rangkap}, \omterm{def:g7:atoms-and-molecules:atom}{atom} yang bermuatan negatif atau
bermuatan parsial $\delta^-$. Sebuah
\emph{situs akseptor elektron}\index{situs akseptor elektron} adalah
tempat yang miskin \omterm{def:g9:inside-the-atom:nucleus}{elektron}, yang mampu menerima sepasang \omterm{def:g9:inside-the-atom:nucleus}{elektron}: \omterm{def:g7:atoms-and-molecules:atom}{atom}
yang bermuatan positif atau bermuatan parsial $\delta^+$.
\end{definition}

\begin{example}[Situs pada dua molekul]\label{ex:g12:curly-arrows:sites}
Dalam bromoetana, \ce{CH3-CH2-Br}, brom lebih elektronegatif daripada
karbon: karbon yang terikat padanya bermuatan $\delta^+$ (situs
akseptor), bromnya $\delta^-$ dan memiliki tiga \omterm{def:g10:lewis-and-shape:lone-pair}{pasangan elektron bebas}.
Dalam propanon, oksigen gugus \ce{C=O} bermuatan $\delta^-$ dan memiliki
dua \omterm{def:g10:lewis-and-shape:lone-pair}{pasangan elektron bebas} (situs donor), karbonnya $\delta^+$ (situs
akseptor). \omterm{def:g9:ions:ion}{Ion} hidroksida \ce{OH-}, dengan muatan negatif dan tiga
\omterm{def:g10:lewis-and-shape:lone-pair}{pasangan elektron bebasnya}, adalah donor yang kuat.
\end{example}

\begin{omfigure}
\setchemfig{atom sep=2.2em}
\begin{tabular}{cc}
\chemfig{H_3C-C(-[2]H)(-[6]H)-\charge{90:2pt=\:,0=\:,270:2pt=\:}{Br}} &
\chemfig{H_3C-C(-[6]CH_3)=[2]\charge{0=\:,180=\:}{O}} \\[2pt]
\small $\delta^+$ pada karbon yang terikat pada \ce{Br}, $\delta^-$ pada \ce{Br} &
\small $\delta^+$ pada karbon \ce{C=O}, $\delta^-$ pada \ce{O} \\
\small bromoetana & \small propanon
\end{tabular}

{\small \omterm{def:g7:atoms-and-molecules:atom}{Atom-atom} karbon yang miskin \omterm{def:g9:inside-the-atom:nucleus}{elektron} ($\delta^+$, situs
akseptor) di sebelah \omterm{def:g7:atoms-and-molecules:atom}{atom-atom} elektronegatif yang kaya \omterm{def:g9:inside-the-atom:nucleus}{elektron}
($\delta^-$, dengan \omterm{def:g10:lewis-and-shape:lone-pair}{pasangan elektron bebas}: situs donor).}
\end{omfigure}

\section{Panah lengkung}

\begin{definition}[Panah lengkung]\label{def:g12:curly-arrows:curly-arrow}
Sebuah \emph{panah lengkung}\index{panah lengkung} memperlihatkan
perpindahan sepasang \omterm{def:g9:inside-the-atom:nucleus}{elektron} selama satu tahap reaksi. Panah itu
berawal dari pasangan yang berpindah, yaitu \omterm{def:g10:lewis-and-shape:lone-pair}{pasangan elektron bebas} atau
sebuah ikatan, dan berakhir di tempat pasangan itu pergi: pada sebuah
\omterm{def:g7:atoms-and-molecules:atom}{atom}, tempat pasangan itu membentuk ikatan baru, atau pada \omterm{def:g7:atoms-and-molecules:atom}{atom} yang
mengambil pasangan \omterm{def:g9:inside-the-atom:nucleus}{elektron} sebuah ikatan yang putus.
\end{definition}

\begin{method}[Menggambar panah lengkung]\label{met:g12:curly-arrows:drawing}
\begin{enumerate}
  \item Temukan situs donor dan situs akseptornya.
  \item Gambarlah panah dari pasangan donor (\omterm{def:g10:lewis-and-shape:lone-pair}{pasangan elektron bebas} atau
        sebuah ikatan) ke \omterm{def:g7:atoms-and-molecules:atom}{atom} akseptor: di situlah ikatan baru
        terbentuk.
  \item Jika \omterm{def:g7:atoms-and-molecules:atom}{atom} akseptor itu lalu akan memiliki terlalu banyak \omterm{def:g9:inside-the-atom:nucleus}{elektron}
        (lebih dari satu oktet, atau lebih dari dua untuk hidrogen),
        gambarlah panah kedua yang memutus salah satu ikatannya, dengan
        pasangannya pergi ke \omterm{def:g7:atoms-and-molecules:atom}{atom} yang lebih elektronegatif.
  \item Periksalah muatan sesudah tahap itu: \omterm{def:g7:atoms-and-molecules:atom}{atom} yang memberikan
        \omterm{def:g10:lewis-and-shape:lone-pair}{pasangan elektron bebas} menjadi satu satuan lebih positif, \omterm{def:g7:atoms-and-molecules:atom}{atom}
        yang menerima pasangan sebagai \omterm{def:g10:lewis-and-shape:lone-pair}{pasangan elektron bebas} menjadi
        satu satuan lebih negatif; muatan totalnya tidak berubah.
\end{enumerate}
\end{method}

\begin{proposition}[Muatan kekal dalam setiap tahap]\label{prop:g12:curly-arrows:charge}
Muatan listrik total spesies-spesies yang terlibat sama sebelum dan
sesudah setiap tahap sebuah mekanisme.
\end{proposition}

\begin{proof}
\omterm{def:g12:curly-arrows:curly-arrow}{Panah lengkung} hanya memindahkan \omterm{def:g9:inside-the-atom:nucleus}{elektron} yang sudah ada; tidak ada
\omterm{def:g9:inside-the-atom:nucleus}{elektron} yang diciptakan atau dimusnahkan, dan inti-inti \omterm{def:g7:atoms-and-molecules:atom}{atom} tidak
berubah.
\end{proof}

\begin{example}[Panah pertama]\label{ex:g12:curly-arrows:first}
Sebuah \omterm{def:g9:ions:ion}{ion} hidrogen \ce{H+} bertemu dengan \omterm{def:g9:ions:ion}{ion} hidroksida \ce{OH-}.
Sepasang \omterm{def:g9:inside-the-atom:nucleus}{elektron} bebas oksigen membentuk ikatan dengan \omterm{def:g9:ions:ion}{ion} hidrogen,
yang tidak memiliki \omterm{def:g9:inside-the-atom:nucleus}{elektron}: \ce{H+ + OH- -> H2O}. Satu panah, dari
\omterm{def:g10:lewis-and-shape:lone-pair}{pasangan elektron bebas} oksigen ke \ce{H+}; muatan $+1$ dan $-1$ menjadi
0 dalam \omterm{def:g7:atoms-and-molecules:molecule}{molekul} air.
\end{example}

\begin{omfigure}
\setchemfig{atom sep=2.1em}
\chemfig{@{hp}H^{+}}\qquad\raisebox{0.2ex}{$+$}\qquad
\chemfig{H-@{ox}\charge{90=\:,0=\:,270=\:}{O}^{-}}
\qquad\raisebox{0.2ex}{$\longrightarrow$}\qquad
\chemfig{H-\charge{90=\:,270=\:}{O}-H}
\chemmove{
  \draw[->, thick, omProp, shorten <=4pt, shorten >=3pt]
    ($(ox)+(0,0.25)$) .. controls +(110:9mm) and +(60:9mm) .. (hp);}

{\small Tahap yang paling sederhana: sepasang \omterm{def:g9:inside-the-atom:nucleus}{elektron} bebas \omterm{def:g9:ions:ion}{ion}
hidroksida berpindah dan membentuk ikatan \ce{O-H} yang baru dengan \omterm{def:g9:ions:ion}{ion}
hidrogen.}
\end{omfigure}

\section{Tahap elementer dan zat antara}

\begin{definition}[Mekanisme, tahap elementer, zat antara]\label{def:g12:curly-arrows:mechanism}
Yang disebut \emph{mekanisme reaksi}\index{mekanisme reaksi} sebuah
reaksi adalah rangkaian tahap yang mengubah \omterm{def:g7:chemical-reactions:reactant}{pereaksi} menjadi produk.
Setiap \emph{tahap elementer}\index{tahap elementer} adalah satu
peristiwa tunggal tempat beberapa pasang \omterm{def:g9:inside-the-atom:nucleus}{elektron} berpindah sekaligus.
Spesies yang terbentuk dalam satu tahap dan habis terpakai dalam tahap
berikutnya disebut \emph{zat antara reaksi}\index{zat antara reaksi}:
spesies itu tidak muncul dalam persamaan keseluruhan.
\end{definition}


\begin{example}[Zat antara karbokation]\label{ex:g12:curly-arrows:carbocation}
Ketika hidrogen bromida diadisikan pada etena, reaksinya berlangsung
dalam dua tahap. Pada tahap pertama, \omterm{def:g10:lewis-and-shape:multiple-bond}{ikatan rangkap} etena memberikan
sepasang \omterm{def:g9:inside-the-atom:nucleus}{elektron} kepada hidrogen \ce{HBr}, dan pasangan \ce{H-Br} pergi
ke brom, yang lepas sebagai \ce{Br-}; karbon yang kehilangan bagiannya
dari \omterm{def:g10:lewis-and-shape:multiple-bond}{ikatan rangkap} tinggal memiliki enam \omterm{def:g9:inside-the-atom:nucleus}{elektron} saja dan bermuatan
positif: sebuah karbokation, \ce{CH3-CH2+}. Pada tahap kedua, sepasang
\omterm{def:g9:inside-the-atom:nucleus}{elektron} bebas \ce{Br-} membentuk ikatan dengan karbon itu. Karbokation
itu adalah zat antara.
\end{example}

\section{Tiga kategori reaksi}

\begin{definition}[Substitusi, adisi, eliminasi]\label{def:g12:curly-arrows:categories}
Dalam \emph{substitusi}\index{substitusi}, sebuah \omterm{def:g7:atoms-and-molecules:atom}{atom} atau gugus pada
\omterm{def:g7:atoms-and-molecules:molecule}{molekul} digantikan oleh \omterm{def:g7:atoms-and-molecules:atom}{atom} atau gugus lain. Dalam
\emph{adisi}\index{adisi}, dua \omterm{def:g7:atoms-and-molecules:molecule}{molekul} bergabung menjadi satu, dan
sebuah \omterm{def:g10:lewis-and-shape:multiple-bond}{ikatan rangkap} menjadi ikatan tunggal. Dalam
\emph{eliminasi}\index{eliminasi}, sebuah \omterm{def:g7:atoms-and-molecules:molecule}{molekul} kecil (sering kali
air) dilepaskan dari sebuah \omterm{def:g7:atoms-and-molecules:molecule}{molekul}, dan sebuah ikatan tunggal menjadi
\omterm{def:g10:lewis-and-shape:multiple-bond}{ikatan rangkap}.
\end{definition}

\begin{omfigure}
\setchemfig{atom sep=2em}
\small
\textbf{\omterm{def:g12:curly-arrows:categories}{Substitusi}} (satu tahap): \ce{CH3-CH2-Br + OH- -> CH3-CH2-OH + Br-}

\medskip
\chemfig{H-@{o1}\charge{90=\:,0=\:,270=\:}{O}^{-}}\qquad
\chemfig{H_3C-@{c1}C(-[2]H)(-[6]H)-[@{b1}]@{br1}\charge{90=\:,0=\:,270=\:}{Br}}
\qquad$\longrightarrow$\qquad
\chemfig{H_3C-C(-[2]H)(-[6]H)-O-H}\quad$+$\quad\ce{Br-}
\chemmove{
  \draw[->, thick, omProp, shorten <=4pt, shorten >=3pt]
    ($(o1)+(0.15,0.2)$) .. controls +(60:9mm) and +(120:9mm) .. (c1);
  \draw[->, thick, omDef, shorten <=2pt, shorten >=4pt]
    (b1) .. controls +(-90:7mm) and +(-90:7mm) .. (br1);}

\vspace{6mm}
\textbf{\omterm{def:g12:curly-arrows:categories}{Adisi}} (dua tahap): \ce{CH2=CH2 + HBr -> CH3-CH2-Br}

\vspace{9mm}
\chemfig{H_2C=[@{pi}]CH_2}\quad$+$\quad
\chemfig{@{h2}H-[@{b2}]@{br2}Br}
\qquad$\longrightarrow$\qquad
\chemfig{H_3C-@{cp}\charge{90:2pt=$\scriptstyle +$}{C}H_2}\quad$+$\quad
\chemfig{@{brm}\charge{90=\:,180=\:,270=\:,0=\:}{Br}^{-}}
\qquad$\longrightarrow$\qquad \ce{CH3-CH2-Br}
\chemmove{
  \draw[->, thick, omProp, shorten <=2pt, shorten >=3pt]
    (pi) .. controls +(90:10mm) and +(110:9mm) .. (h2);
  \draw[->, thick, omDef, shorten <=2pt, shorten >=4pt]
    (b2) .. controls +(-90:7mm) and +(-90:7mm) .. (br2);
  \draw[->, thick, omProp, shorten <=4pt, shorten >=3pt]
    ($(brm)+(0,-0.25)$) .. controls +(-120:8mm) and +(-60:8mm) .. (cp);}

\vspace{6mm}
\textbf{\omterm{def:g12:curly-arrows:categories}{Eliminasi}} (dehidrasi etanol dengan asam, tiga tahap):
\ce{CH3-CH2-OH -> CH2=CH2 + H2O}

\medskip
\begin{tabular}{@{}l@{}}
1.\ oksigen mengambil \omterm{def:g9:ions:ion}{ion} hidrogen: \ce{CH3-CH2-OH + H+ -> CH3-CH2-OH2+};\\
2.\ pasangan \ce{C-O} lepas bersama air: \ce{CH3-CH2-OH2+ -> CH3-CH2+ + H2O};\\
3.\ sepasang \omterm{def:g9:inside-the-atom:nucleus}{elektron} \ce{C-H} tetangga menjadi \omterm{def:g10:lewis-and-shape:multiple-bond}{ikatan rangkap}, dan
    \ce{H+} lepas:\\
\phantom{3.\ }\ce{CH3-CH2+ -> CH2=CH2 + H+}.
\end{tabular}

\medskip
{\small Tiga mekanisme. Panah jingga: sepasang \omterm{def:g9:inside-the-atom:nucleus}{elektron} donor membentuk
ikatan baru. Panah biru: sebuah ikatan putus, pasangannya pergi ke \omterm{def:g7:atoms-and-molecules:atom}{atom}
yang lebih elektronegatif. Dalam \omterm{def:g12:curly-arrows:categories}{eliminasi}, \omterm{def:g9:ions:ion}{ion} hidrogen dikembalikan
pada akhirnya.}
\end{omfigure}


\begin{remark}[Mengubah kerangka atau gugus]\label{rem:g12:curly-arrows:skeleton}
Sebagian reaksi hanya mengubah \omterm{def:g11:functional-groups:functional-group}{gugus fungsi} dan mempertahankan kerangka
karbonnya: \omterm{def:g12:curly-arrows:categories}{substitusi} bromoetana mengubah \omterm{def:g11:functional-groups:halogenoalkane}{haloalkana} menjadi \omterm{def:g11:functional-groups:alcohol}{alkohol},
dua karbon sebelumnya, dua karbon sesudahnya. Reaksi lain membangun atau
memotong kerangka: membentuk ikatan karbon--karbon yang baru
memperpanjang rantai, tahap kunci dalam membuat \omterm{def:g7:atoms-and-molecules:molecule}{molekul} besar dari
\omterm{def:g7:atoms-and-molecules:molecule}{molekul-molekul} kecil. Merencanakan \omterm{def:g10:synthesis-yield:synthesis}{sintesis} berarti memilih reaksi dari
kedua jenis itu, dengan urutan yang tepat.
\end{remark}

\section{Latihan}

\begin{exercise}[$\star$]\label{exo:g12:curly-arrows:1}
Pada setiap spesies, sebutkan sebuah situs donor dan, jika ada, sebuah
situs akseptor: \ce{OH-}; \ce{H+}; \ce{CH2=CH2}; \ce{CH3-Cl}; \ce{H2O}.
\end{exercise}

\begin{exercise}[$\star$]\label{exo:g12:curly-arrows:2}
\omterm{def:g12:curly-arrows:categories}{Substitusi}, \omterm{def:g12:curly-arrows:categories}{adisi}, atau \omterm{def:g12:curly-arrows:categories}{eliminasi}?
(a) \ce{CH3-CH2-Cl + OH- -> CH3-CH2-OH + Cl-};
(b) \ce{CH2=CH2 + H2O -> CH3-CH2-OH};
(c) \ce{CH3-CH2-OH -> CH2=CH2 + H2O};
(d) \ce{CH2=CH2 + Br2 -> CH2Br-CH2Br}.
\end{exercise}

\begin{exercise}[$\star$]\label{exo:g12:curly-arrows:3}
Dari mana \omterm{def:g12:curly-arrows:curly-arrow}{panah lengkung} berawal, dan di mana panah itu berakhir? Apa
arti \omterm{def:g12:curly-arrows:curly-arrow}{panah lengkung} dari sepasang \omterm{def:g9:inside-the-atom:nucleus}{elektron} bebas \omterm{def:g7:atoms-and-molecules:atom}{atom} oksigen ke sebuah
\omterm{def:g7:atoms-and-molecules:atom}{atom} karbon?
\end{exercise}

\begin{exercise}[$\star$]\label{exo:g12:curly-arrows:4}
Dalam reaksi \ce{H+ + NH3 -> NH4+}, spesies mana yang memberikan
pasangan \omterm{def:g9:inside-the-atom:nucleus}{elektronnya}? Gambarlah panahnya. Berapa muatan produknya?
\end{exercise}

\begin{exercise}[$\star$]\label{exo:g12:curly-arrows:5}
Apa itu \omterm{def:g12:curly-arrows:mechanism}{zat antara reaksi}? Sebutkan zat antara pada \omterm{def:g12:curly-arrows:categories}{adisi} \ce{HBr}
pada etena.
\end{exercise}

\begin{exercise}[$\star\star$]\label{exo:g12:curly-arrows:6}
Gambarlah panah-panah lengkung \omterm{def:g12:curly-arrows:categories}{substitusi}
\ce{CH3-Br + OH- -> CH3-OH + Br-} (satu tahap), lalu periksalah
muatannya.
\end{exercise}

\begin{exercise}[$\star\star$]\label{exo:g12:curly-arrows:7}
Pada tahap pertama \omterm{def:g12:curly-arrows:categories}{adisi} \ce{HBr} pada etena, mengapa pasangan \ce{H-Br}
pergi ke brom dan bukan ke hidrogen? Berapa muatan setiap produk tahap
itu?
\end{exercise}

\begin{exercise}[$\star\star$]\label{exo:g12:curly-arrows:8}
Pada dehidrasi etanol, tuliskan ketiga tahapnya dan sebutkan, untuk
setiap tahap, spesies mana yang menjadi donor dan mana yang menjadi
akseptor. Mengapa \omterm{def:g9:ions:ion}{ion} hidrogen tidak dituliskan dalam persamaan
keseluruhan?
\end{exercise}

\begin{exercise}[$\star\star$]\label{exo:g12:curly-arrows:9}
Tahap \ce{CH3-CH2+ + H2O -> CH3-CH2-OH2+} menyusul \omterm{def:g12:curly-arrows:categories}{adisi} \omterm{def:g9:ions:ion}{ion} hidrogen
pada etena di dalam air. Gambarlah panahnya. Apa yang terbentuk setelah
\ce{H+} lepas lagi? Tuliskan persamaan keseluruhannya dan tentukan
kategorinya.
\end{exercise}

\begin{exercise}[$\star\star$]\label{exo:g12:curly-arrows:10}
Bacalah mekanisme \omterm{def:g12:curly-arrows:categories}{substitusi} bromoetana: berapa pasang \omterm{def:g9:inside-the-atom:nucleus}{elektron} yang
berpindah? Ikatan mana yang terbentuk, mana yang putus? Adakah zat
antara?
\end{exercise}

\begin{exercise}[$\star\star$]\label{exo:g12:curly-arrows:11}
% ledger: en:C, en:Cl, en:Br
Klorometana bereaksi dengan \omterm{def:g9:ions:ion}{ion} hidroksida lebih lambat daripada
bromometana. Dengan \omterm{def:g11:polarity-and-cohesion:electronegativity}{keelektronegatifan} (C 2.55, Cl 3.16, Br 2.96),
jelaskan karbon mana yang lebih $\delta^+$. Apakah hal itu menjelaskan
kecepatannya? (Pikirkan juga seberapa mudah \omterm{def:g9:ions:ion}{ion} yang lepas mengambil
pasangan \omterm{def:g9:inside-the-atom:nucleus}{elektronnya}.)
\end{exercise}

\begin{exercise}[$\star\star\star$]\label{exo:g12:curly-arrows:12}
Hidrogen klorida diadisikan pada propena, \ce{CH2=CH-CH3}. Pada tahap
pertama, \omterm{def:g9:ions:ion}{ion} hidrogen dapat berikatan dengan karbon mana pun dari \omterm{def:g10:lewis-and-shape:multiple-bond}{ikatan
rangkap} itu, sehingga ada dua karbokation yang mungkin. Gambarlah
keduanya, kedua produk yang mungkin, lalu namailah. (Produk mana yang
paling banyak terbentuk dipelajari dalam buku Tahun Pertama
universitas.)
\end{exercise}

\begin{exercise}[$\star\star\star$]\label{exo:g12:curly-arrows:13}
Seorang siswa menggambar tahap pertama \omterm{def:g12:curly-arrows:categories}{substitusi} bromoetana dengan
panah dari \omterm{def:g7:atoms-and-molecules:atom}{atom} brom ke oksigen \ce{OH-}. Jelaskan kedua kesalahannya.
\end{exercise}

\begin{exercise}[$\star\star\star$]\label{exo:g12:curly-arrows:14}
Dalam pembentukan \omterm{def:g11:functional-groups:carboxylic-acid}{ester} dari \omterm{def:g11:functional-groups:carboxylic-acid}{asam karboksilat} dan \omterm{def:g11:functional-groups:alcohol}{alkohol}, oksigen
\omterm{def:g11:functional-groups:alcohol}{alkohol} menyerang karbon gugus \ce{C=O}. Dengan \omterm{def:g11:polarity-and-cohesion:polar-bond}{muatan parsial},
jelaskan mengapa karbon itu merupakan situs akseptor dan oksigen itu
situs donor. \omterm{def:g7:atoms-and-molecules:molecule}{Molekul} kecil apa yang dilepaskan secara keseluruhan, dan
termasuk kategori apa reaksi keseluruhannya?
\end{exercise}

\begin{exercise}[$\star\star\star$]\label{exo:g12:curly-arrows:15}
Etena bereaksi dengan brom \ce{Br2}, sebuah \omterm{def:g11:polarity-and-cohesion:polar-molecule}{molekul nonpolar}. Ketika
sebuah \omterm{def:g7:atoms-and-molecules:molecule}{molekul} \ce{Br2} mendekati \omterm{def:g10:lewis-and-shape:multiple-bond}{ikatan rangkap}, \omterm{def:g9:inside-the-atom:nucleus}{elektron-elektron}
\ce{Br-Br} terdorong menjauh dari \omterm{def:g7:atoms-and-molecules:atom}{atom} brom yang dekat. Jelaskan
mengapa hal ini menciptakan sebuah situs akseptor, lalu usulkan tahap
pertama mekanismenya.
\end{exercise}

\section{Soal: Membuat Bromoetana}

\begin{problem}[{Soal akhir pekan --- dari sepuluh gram etena: berapa
bromoetana?}]
\label{pb:g12:curly-arrows:1}
% ledger: aw:C, aw:H, aw:Br, en:H, en:Br, en:C
Bromoetana, bahan awal bagi banyak \omterm{def:g10:synthesis-yield:synthesis}{sintesis}, dibuat dengan
mengadisikan hidrogen bromida pada etena:
\ce{CH2=CH2 + HBr -> CH3-CH2-Br}. Seorang ahli kimia memakai
\qty{10.0}{g} etena dan hidrogen bromida berlebih; \omterm{def:g10:synthesis-yield:yield}{rendemennya}
\qty{75}{\percent}.

\medskip
\noindent\textbf{Bagian I --- Situs-situsnya.}
\begin{enumerate}
  \item Gambarlah \omterm{def:g10:lewis-and-shape:lewis-structure}{struktur Lewis} etena dan hidrogen bromida.
  \item Di mana situs donor etena? Mengapa?
  \item Hitung selisih \omterm{def:g11:polarity-and-cohesion:electronegativity}{keelektronegatifan} ikatan \ce{H-Br}
        (H 2.20, Br 2.96). \omterm{def:g7:atoms-and-molecules:atom}{Atom} mana yang bermuatan $\delta^+$?
  \item \omterm{def:g7:atoms-and-molecules:atom}{Atom} \ce{HBr} mana yang menjadi situs akseptor?
\end{enumerate}

\medskip
\noindent\textbf{Bagian II --- Mekanismenya.}
\begin{enumerate}[resume]
  \item Gambarlah tahap pertama dengan kedua \omterm{def:g12:curly-arrows:curly-arrow}{panah lengkungnya}.
  \item Sebutkan kedua spesies yang terbentuk beserta muatannya.
  \item Periksalah kekekalan muatan pada tahap pertama.
  \item Gambarlah tahap kedua dengan \omterm{def:g12:curly-arrows:curly-arrow}{panah lengkungnya}.
  \item Berapa \omterm{def:g9:inside-the-atom:nucleus}{elektron} yang mengelilingi karbon positif zat antara itu?
        Mengapa karbon itu begitu reaktif?
\end{enumerate}

\medskip
\noindent\textbf{Bagian III --- Reaksi jenis apa?}
\begin{enumerate}[resume]
  \item Apa zat antara mekanisme ini?
  \item Mengapa zat antara itu tidak muncul dalam persamaan
        keseluruhan?
  \item Termasuk kategori apa reaksi ini?
  \item Apakah reaksi ini mengubah kerangka karbon, atau hanya \omterm{def:g11:functional-groups:functional-group}{gugus
        fungsinya}?
\end{enumerate}

\medskip
\noindent\textbf{Bagian IV --- Berapa banyak produknya?}
\begin{enumerate}[resume]
  \item Hitung \omterm{def:g10:the-mole:molar-mass}{massa molar} etena dan bromoetana.
  \item Hitung jumlah etena.
  \item Isilah \omterm{def:g10:reaction-progress-table:extent}{tabel kemajuannya}. \omterm{def:g7:chemical-reactions:reactant}{Pereaksi} mana yang menjadi pembatas?
  \item Hitung massa maksimum bromoetana.
  \item Hitung massa yang diperoleh dengan \omterm{def:g10:synthesis-yield:yield}{rendemen} \qty{75}{\percent}.
  \item Nyatakan jawaban akhirnya: berapa massa bromoetana yang diperoleh
        ahli kimia itu?
\end{enumerate}
\end{problem}
